% Lecture 8 notes: \input by notes.tex, \subimport'ed by the semester book.
% Exam questions (экзамен): 18 (уравнение Вейерштрасса, определение, дискриминант),
% 19 (сложение точек, группа точек), 20 (кривые над конечными полями, теорема Хассе).
\begin{lecture}{8}{Эллиптические кривые}

\section{Зачем криптографии эллиптические кривые}

Протоколы Диффи"--~Хеллмана и Эль-Гамаля (лекция~1) работают в любой конечной циклической
группе: нужно лишь уметь быстро умножать элементы и чтобы задача дискретного логарифмирования
(DLP) в группе была трудной. В группе $\Znx{p}$ DLP решается субэкспоненциальным методом
исчисления индексов (лекция~4), поэтому модуль $p$ приходится брать длиной $3072$ бита и больше.
В 1985 г. Н.~Коблиц~\cite{koblitz-ecc} и В.~Миллер~\cite{miller-ecc} независимо предложили
использовать \emph{группу точек эллиптической кривой} над конечным полем. Для неё (кроме
специальных классов кривых) не известно ничего лучше универсальных алгоритмов сложности
$\bigO(\sqrt n)$ (лекции~3 и~9), и ключ длиной $256$ бит даёт ту же стойкость (около $128$ бит),
что $3072$-битный модуль RSA или Диффи"--~Хеллмана~\cite{nist-sp800-57}. Сегодня
эллиптическая криптография "--- стандарт де-факто: подпись ECDSA и Ed25519, обмен ключами ECDH
в TLS, кривая secp256k1 в Bitcoin.

Кроме того, эллиптические кривые дают один из лучших алгоритмов факторизации "--- метод Ленстры
(лекция~10) "--- и используются в доказательстве простоты больших чисел.

В этой лекции мы определим эллиптические кривые, введём на множестве их точек структуру
абелевой группы и изучим, сколько точек имеет кривая над конечным полем $\GF{p}$.
Всюду $K$ "--- поле (чаще всего $\RR$, $\QQ$ или $\GF{p}$), $\overline K$ "--- его алгебраическое
замыкание.

\section{Уравнение Вейерштрасса}

\subsection{Общая и короткая формы}

\begin{definition}
\emph{Уравнением Вейерштрасса} над полем $K$ называется уравнение
\begin{equation}\label{l8:general}
  y^2 + a_1 xy + a_3 y = x^3 + a_2 x^2 + a_4 x + a_6, \qquad a_i \in K.
\end{equation}
\end{definition}

Если $\operatorname{char} K \ne 2$, выделим полный квадрат по $y$: замена
$y \mapsto y - \frac{a_1 x + a_3}{2}$ уничтожает слагаемые $a_1 xy$ и $a_3 y$, и уравнение
принимает вид $y^2 = x^3 + b_2 x^2 + b_4 x + b_6$. Если ещё и $\operatorname{char} K \ne 3$,
замена $x \mapsto x - \frac{b_2}{3}$ (выделение полного куба) уничтожает слагаемое с $x^2$.
Обе замены обратимы и переводят многочлены в многочлены, поэтому все свойства, о которых пойдёт
речь, при них сохраняются. Итак, при $\operatorname{char} K \ne 2, 3$ достаточно изучать
\emph{короткую форму Вейерштрасса}
\begin{equation}\label{l8:short}
  E\colon\ y^2 = f(x) = x^3 + ax + b, \qquad a, b \in K.
\end{equation}
Всюду далее, кроме раздела~\ref{l8:char23}, предполагается $\operatorname{char} K \ne 2, 3$
(для криптографии над $\GF{p}$ это значит $p > 3$).

\subsection{Особые точки и дискриминант}

Положим $F(x, y) = y^2 - f(x)$, так что кривая задаётся уравнением $F(x, y) = 0$.

\begin{definition}
Точка $(x_0, y_0)$ кривой $F = 0$ называется \emph{особой}, если в ней обе частные производные
обращаются в нуль: $F_x(x_0, y_0) = F_y(x_0, y_0) = 0$. Кривая без особых точек (над
$\overline K$) называется \emph{гладкой}, или \emph{неособой}.
\end{definition}

В неособой точке у кривой есть единственная касательная
$F_x(x_0, y_0)(x - x_0) + F_y(x_0, y_0)(y - y_0) = 0$, а в особой точке касательная не
определена: кривая имеет там «самопересечение» или «клюв» (рис.~ниже). Именно касательные
понадобятся нам для удвоения точек.

\begin{definition}
\emph{Дискриминантом} кривой~\eqref{l8:short} называется $\Delta = -16(4a^3 + 27b^2)$.
\emph{Эллиптической кривой} над $K$ называется кривая~\eqref{l8:short} с $\Delta \ne 0$ вместе
с бесконечно удалённой точкой $\ECpt$ (раздел~\ref{l8:infinity}).
\end{definition}

\begin{theorem}\label{l8:disc}
Пусть $\operatorname{char} K \ne 2, 3$. Следующие условия равносильны:
\begin{enumerate}
  \item кривая $y^2 = x^3 + ax + b$ не имеет особых точек над $\overline K$;
  \item многочлен $f(x) = x^3 + ax + b$ не имеет кратных корней в $\overline K$;
  \item $4a^3 + 27b^2 \ne 0$, т.\,е. $\Delta \ne 0$.
\end{enumerate}
Если же $4a^3 + 27b^2 = 0$, то кратный корень равен $x_0 = -\frac{3b}{2a}$ при $a \ne 0$
и $x_0 = 0$ при $a = 0$, а единственная особая точка "--- $(x_0, 0)$.
\end{theorem}

\begin{proof}
$(1 \Leftrightarrow 2)$. Частные производные: $F_y = 2y$, $F_x = -f'(x)$. Точка $(x_0, y_0)$
кривой особая тогда и только тогда, когда $2y_0 = 0$, $f'(x_0) = 0$ и $y_0^2 = f(x_0)$, т.\,е.
(так как $2 \ne 0$ в $K$) $y_0 = 0$ и $f(x_0) = f'(x_0) = 0$. Последнее и означает, что $x_0$
"--- кратный корень $f$.

$(2 \Leftrightarrow 3)$. Пусть $x_0$ "--- кратный корень: $f'(x_0) = 3x_0^2 + a = 0$ и
$f(x_0) = x_0(x_0^2 + a) + b = 0$. Из первого равенства $x_0^2 = -\frac a3$, и второе даёт
$\frac{2a}{3} x_0 + b = 0$. Если $a = 0$, то $b = 0$ и $4a^3 + 27b^2 = 0$. Если $a \ne 0$, то
$x_0 = -\frac{3b}{2a}$, и подстановка в $x_0^2 = -\frac a3$ даёт $\frac{9b^2}{4a^2} = -\frac a3$,
т.\,е. $27b^2 = -4a^3$.

Обратно, пусть $4a^3 + 27b^2 = 0$. При $a = 0$ имеем $b = 0$ и $f = x^3$ с тройным корнем $0$.
При $a \ne 0$ положим $x_0 = -\frac{3b}{2a}$; тогда $x_0^2 = \frac{9b^2}{4a^2} = -\frac{a}{3}$
(в силу $27b^2 = -4a^3$), поэтому $f'(x_0) = 3x_0^2 + a = 0$ и
$f(x_0) = x_0(x_0^2 + a) + b = \frac{2a}{3}\cdot\bigl(-\frac{3b}{2a}\bigr) + b = 0$.
\end{proof}

\begin{remark}
Множитель $-16$ в определении $\Delta$ появляется из общей формулы дискриминанта
уравнения~\eqref{l8:general}; для условия $\Delta \ne 0$ он не важен (при $\operatorname{char} K \ne 2$).
Если $e_1, e_2, e_3 \in \overline K$ "--- корни $f$, то $e_1 + e_2 + e_3 = 0$ и
\[
  -4a^3 - 27b^2 = \bigl((e_1 - e_2)(e_1 - e_3)(e_2 - e_3)\bigr)^2, \qquad
  \Delta = 16 \prod_{i < j} (e_i - e_j)^2,
\]
что проверяется раскрытием скобок с $a = e_1 e_2 + e_1 e_3 + e_2 e_3$, $b = -e_1 e_2 e_3$; это
даёт второе доказательство равносильности $(2 \Leftrightarrow 3)$. Над $\RR$ отсюда следует:
$\Delta > 0$ тогда и только тогда, когда у $f$ три различных вещественных корня (если $e_2$
и $e_3 = \overline{e_2}$ комплексно сопряжены, то $(e_2 - e_3)^2 < 0$, а
$(e_1 - e_2)^2 (e_1 - e_3)^2 = \abs{e_1 - e_2}^4 > 0$).
\end{remark}

Вещественная кривая при $\Delta > 0$ состоит из двух компонент («овал» и бесконечная ветвь),
при $\Delta < 0$ "--- из одной. При $\Delta = 0$ кривая имеет особую точку: \emph{узел}
(двойной корень, $f = (x - 1)^2(x + 2)$) или \emph{точку возврата} (тройной корень, $f = x^3$).

\begin{center}
\begin{tikzpicture}[scale=0.62]
  \begin{scope}
    \draw[-{Stealth[length=1.5mm]},caGray] (-1.8,0) -- (2.4,0); \draw[-{Stealth[length=1.5mm]},caGray] (0,-2.4) -- (0,2.4);
    \draw[caBlue,thick,smooth,samples=60,domain=-1:1] plot ({-1+\x*\x}, {\x*sqrt(max(0,(-1+\x*\x)*(-2+\x*\x)))});
    \draw[caBlue,thick,smooth,samples=60,domain=-1.15:1.15] plot ({1+\x*\x}, {\x*sqrt((1+\x*\x)*(2+\x*\x))});
    \node[below] at (0.3,-2.5) {\small $y^2 = x^3 - x$, $\Delta > 0$};
  \end{scope}
  \begin{scope}[xshift=7.8cm]
    \draw[-{Stealth[length=1.5mm]},caGray] (-1.8,0) -- (2.4,0); \draw[-{Stealth[length=1.5mm]},caGray] (0,-2.4) -- (0,2.4);
    \draw[caBlue,thick,smooth,samples=80,domain=-1.5:1.5] plot ({-1+\x*\x}, {\x*sqrt((-1+\x*\x)^2-(-1+\x*\x)+1)});
    \node[below] at (0.3,-2.5) {\small $y^2 = x^3 + 1$, $\Delta < 0$};
  \end{scope}
  \begin{scope}[yshift=-6.2cm]
    \draw[-{Stealth[length=1.5mm]},caGray] (-2.3,0) -- (2.4,0); \draw[-{Stealth[length=1.5mm]},caGray] (0,-2.4) -- (0,2.4);
    \draw[caRed,thick,smooth,samples=80,domain=-1.95:1.95] plot ({\x*\x-2}, {0.55*\x*(\x*\x-3)});
    \fill[caRed] (1,0) circle (2.5pt);
    \node[below] at (0.1,-2.5) {\small $y^2 = (x - 1)^2 (x + 2)$: узел};
  \end{scope}
  \begin{scope}[xshift=7.8cm,yshift=-6.2cm]
    \draw[-{Stealth[length=1.5mm]},caGray] (-1.8,0) -- (2.4,0); \draw[-{Stealth[length=1.5mm]},caGray] (0,-2.4) -- (0,2.4);
    \draw[caRed,thick,smooth,samples=60,domain=-1.3:1.3] plot ({\x*\x}, {\x*\x*\x});
    \fill[caRed] (0,0) circle (2.5pt);
    \node[below] at (0.3,-2.5) {\small $y^2 = x^3$: точка возврата};
  \end{scope}
\end{tikzpicture}
\end{center}
(На рисунке с узлом масштаб по оси $y$ уменьшен почти вдвое.)

Проверка условия $\Delta \ne 0$ "--- первый шаг любой задачи о конкретной кривой:

\taskexample{ec-singular}

\checkyourself{Задаёт ли уравнение $y^2 = x^3 + 10x + 5$ эллиптическую кривую над $\GF{11}$?
А над $\GF{13}$? Найдите особую точку, если она есть.}

\section{Бесконечно удалённая точка}\label{l8:infinity}

Чтобы сложение точек было определено всегда, к аффинной кривой добавляют ещё одну точку.
Её естественно описать на проективной плоскости.

\begin{definition}
\emph{Проективная плоскость} $\mathbb P^2(K)$ "--- множество троек $(X, Y, Z) \in K^3 \setminus
\{(0, 0, 0)\}$ с точностью до пропорциональности: $(X : Y : Z) = (cX : cY : cZ)$ при
$c \in K^*$.
\end{definition}

Точки с $Z \ne 0$ взаимно однозначно соответствуют точкам аффинной плоскости:
$(X : Y : Z) \leftrightarrow (X/Z, Y/Z)$; точки с $Z = 0$ называются \emph{бесконечно
удалёнными}. Подставляя $x = X/Z$, $y = Y/Z$ в~\eqref{l8:short} и умножая на $Z^3$, получаем
однородное уравнение
\[
  Y^2 Z = X^3 + a X Z^2 + b Z^3.
\]
При $Z = 0$ оно превращается в $X^3 = 0$, т.\,е. $X = 0$, и бесконечно удалённая точка на кривой
ровно одна: $\ECpt = (0 : 1 : 0)$. Вертикальная прямая $x = c$, т.\,е. $X = cZ$, проходит через
$\ECpt$; все вертикальные прямые «пересекаются в бесконечности» в точке $\ECpt$. Проверьте
частными производными однородного уравнения, что $\ECpt$ "--- неособая точка, а касательная в
ней "--- бесконечно удалённая прямая $Z = 0$, пересекающая кривую в $\ECpt$ с кратностью~$3$.

Итак, для поля $L \supseteq K$
\[
  E(L) = \{(x, y) \in L^2 : y^2 = x^3 + ax + b\} \cup \{\ECpt\}.
\]

\section{Сложение точек}

\subsection{Правило хорд и касательных}

Ключевое свойство кубической кривой: \emph{прямая пересекает её ровно в трёх точках с учётом
кратности} (включая $\ECpt$). Отсюда «геометрическое» определение суммы:

\begin{quote}
Чтобы сложить $P$ и $Q$, проведём через них прямую (при $P = Q$ "--- касательную); она
пересечёт кривую в третьей точке $R$. Тогда $P + Q$ "--- точка, симметричная $R$ относительно
оси $x$: $P + Q = -R$.
\end{quote}

\begin{center}
\ecaddfigure[scale=0.62]\qquad\qquad\ecdoublefigure[scale=0.62]
\end{center}

Если прямая вертикальна ($Q = (x_1, -y_1)$, или $P = Q$ с $y_1 = 0$), то третья точка "---
$\ECpt$, и $P + Q = \ECpt$. Точка $\ECpt$ служит нулём: прямая через $P$ и $\ECpt$ "---
вертикальная, её третья точка пересечения "--- симметричная $P$ точка, и $P + \ECpt = P$.

\begin{lemma}\label{l8:line}
Пусть $P_1 = (x_1, y_1)$, $P_2 = (x_2, y_2) \in E(K)$, и $L\colon y = \lambda(x - x_1) + y_1$ "---
либо прямая через $P_1$ и $P_2$ при $x_1 \ne x_2$, где $\lambda = \frac{y_2 - y_1}{x_2 - x_1}$,
либо касательная в $P_1 = P_2$ при $y_1 \ne 0$, где $\lambda = \frac{3x_1^2 + a}{2y_1}$.
Тогда абсциссы точек пересечения $L \cap E$ "--- корни (с учётом кратности) многочлена
\[
  g(x) = x^3 + ax + b - \bigl(\lambda(x - x_1) + y_1\bigr)^2
  = (x - x_1)(x - x_2)(x - x_3), \qquad x_3 = \lambda^2 - x_1 - x_2 \in K,
\]
и третья точка пересечения равна $(x_3, \lambda(x_3 - x_1) + y_1)$.
\end{lemma}

\begin{proof}
Точка $(x, y)$ лежит на $L \cap E$ тогда и только тогда, когда $y = \lambda(x - x_1) + y_1$
и $g(x) = 0$. Многочлен $g$ "--- кубический со старшим коэффициентом $1$ и коэффициентом
$-\lambda^2$ при $x^2$; по теореме Виета сумма его корней равна $\lambda^2$.
Число $x_1$ "--- корень $g$, так как $P_1 \in E$. Если $x_1 \ne x_2$, то и $x_2$ "--- корень
($P_2 \in L \cap E$). Если $P_1 = P_2$, то $x_1$ "--- двойной корень: $g'(x_1) = f'(x_1) - 2\lambda y_1
= 3x_1^2 + a - (3x_1^2 + a) = 0$. В обоих случаях третий корень равен
$x_3 = \lambda^2 - x_1 - x_2$ и лежит в $K$.
\end{proof}

Формула для $\lambda$ в случае касательной "--- это неявное дифференцирование
$2y\,y' = 3x^2 + a$. Отражая третью точку, получаем $y_3 = -(\lambda(x_3 - x_1) + y_1)$.

\begin{algorithm}[сложение точек эллиптической кривой]\label{l8:addalg}
\begin{algorithmic}[1]
\Require $P_1, P_2 \in E(K)$, $E\colon y^2 = x^3 + ax + b$.
\Ensure $P_3 = P_1 + P_2$.
\If{$P_1 = \ECpt$} \Return $P_2$ \EndIf
\If{$P_2 = \ECpt$} \Return $P_1$ \EndIf
\State $(x_1, y_1) \gets P_1$, $(x_2, y_2) \gets P_2$
\If{$x_1 = x_2$ и $y_1 = -y_2$} \Return $\ECpt$ \Comment{в том числе $P_1 = P_2$, $y_1 = 0$}
\EndIf
\If{$P_1 \ne P_2$} \State $\lambda \gets (y_2 - y_1)/(x_2 - x_1)$
\Else \State $\lambda \gets (3x_1^2 + a)/(2y_1)$ \EndIf
\State $x_3 \gets \lambda^2 - x_1 - x_2$, \quad $y_3 \gets \lambda(x_1 - x_3) - y_1$
\State \Return $(x_3, y_3)$
\end{algorithmic}
\end{algorithm}

Алгоритм выполняет одно деление (обращение) в поле $K$ и несколько умножений.
Заметим, что точка $-P = (x, -y)$ получается бесплатно, поэтому вычитание стоит столько же,
сколько сложение: $P - Q = P + (-Q)$.

\begin{example}\label{l8:ex-q}
$E\colon y^2 = x^3 - 15x + 18$ над $\QQ$, $\Delta = 76032 > 0$ (корни $f$: $3$ и
$\frac{-3 \pm \sqrt{33}}{2}$), $P = (7, 16)$, $Q = (1, 2)$.
\begin{itemize}
  \item $P + Q$: $\lambda = \frac{2 - 16}{1 - 7} = \frac73$, $x_3 = \frac{49}{9} - 7 - 1 = -\frac{23}{9}$,
    $y_3 = \frac73\bigl(7 + \frac{23}{9}\bigr) - 16 = \frac{170}{27}$, т.\,е.
    $P + Q = \bigl(-\frac{23}{9}, \frac{170}{27}\bigr)$.
  \item $2P$: $\lambda = \frac{3 \cdot 49 - 15}{2 \cdot 16} = \frac{33}{8}$,
    $x_3 = \frac{1089}{64} - 14 = \frac{193}{64}$,
    $y_3 = \frac{33}{8}\bigl(7 - \frac{193}{64}\bigr) - 16 = \frac{223}{512}$.
\end{itemize}
Знаменатели координат быстро растут: над $\QQ$ сложность вычислений с точками увеличивается
с каждым шагом. Над конечным полем такой проблемы нет.
\end{example}

\checkyourself{Почему $x_3$ в лемме~\ref{l8:line} всегда лежит в том же поле $K$, что и
координаты $P_1$, $P_2$, хотя третий корень кубического уравнения, вообще говоря, может не
лежать в $K$?}

\subsection{Группа точек}

\begin{theorem}\label{l8:group}
Множество $E(K)$ с операцией сложения алгоритма~\ref{l8:addalg} является абелевой группой
с нулём $\ECpt$; обратный к $P = (x, y)$ элемент "--- $-P = (x, -y)$.
\end{theorem}

\begin{proof}
\emph{Замкнутость.} По лемме~\ref{l8:line} точка $(x_3, -y_3)$ лежит на $E$ и имеет координаты
из $K$; симметричная ей точка $(x_3, y_3)$ также лежит на $E$, так как уравнение содержит лишь
$y^2$.

\emph{Нейтральный и обратный элементы} "--- по первым строкам алгоритма: $P + \ECpt = P$,
$(x, y) + (x, -y) = \ECpt$. Обратный элемент единствен, так как при $x_1 = x_2$ сумма равна
$\ECpt$ лишь при $y_2 = -y_1$.

\emph{Коммутативность.} В случае $P_1 = P_2$ доказывать нечего. Иначе $\lambda$ и
$x_3 = \lambda^2 - x_1 - x_2$ симметричны по $P_1, P_2$, а
$\lambda(x_2 - x_3) - y_2 = \lambda(x_2 - x_3) - y_1 - \lambda(x_2 - x_1) = \lambda(x_1 - x_3) - y_1$,
так как $y_2 = y_1 + \lambda(x_2 - x_1)$. Геометрически: прямая через $P$ и $Q$ "--- та же, что
через $Q$ и $P$.

\emph{Ассоциативность} $(P + Q) + R = P + (Q + R)$ "--- без доказательства (см.
\cite[разд.~2.4]{washington}, \cite[гл.~III, \S\,3]{silverman-aec}). В «общем положении» это
тождество рациональных функций от координат точек, которое проверяется в системе
компьютерной алгебры; трудность "--- в аккуратном разборе многочисленных частных случаев
($P = Q$, $Q = -R$, \dots).
\end{proof}

Вот проверка ассоциативности в общем положении в Sage (коэффициенты $a, b$ кривой выражены
через две точки $P, Q$, третья точка $R$ "--- общая точка той же кривой); программа печатает
\texttt{True True}:
{\small
\begin{verbatim}
R.<x1,y1,x2,y2,x3> = QQ[]; K = R.fraction_field()
x1, y1, x2, y2, x3 = K.gens()
a = (y1^2 - y2^2 - x1^3 + x2^3)/(x1 - x2); b = y1^2 - x1^3 - a*x1
S.<Y> = K[]; L.<y3> = K.extension(Y^2 - (x3^3 + a*x3 + b))
def add(P, Q):
    (u1, v1), (u2, v2) = P, Q
    lam = (v2 - v1)/(u2 - u1); u3 = lam^2 - u1 - u2
    return (u3, lam*(u1 - u3) - v1)
P, Q, R = (L(x1), L(y1)), (L(x2), L(y2)), (L(x3), y3)
A = add(add(P, Q), R); B = add(P, add(Q, R))
print(A[0] == B[0], A[1] == B[1])
\end{verbatim}
}

\begin{remark}
Правило «$P + Q + R = \ECpt$ для трёх точек одной прямой» объясняет выбор знака: третья точка
$R$ пересечения равна $-(P + Q)$. Можно было бы определить «$P * Q = R$» без отражения, но
такая операция не ассоциативна и не имеет нейтрального элемента.
\end{remark}

\begin{remark}
Поскольку формулы сложения "--- рациональные функции с коэффициентами из простого подполя,
а тождество ассоциативности "--- тождество таких функций, теорема~\ref{l8:group} верна над
\emph{любым} полем характеристики $\ne 2, 3$ (а для общей формы~\eqref{l8:general} "--- над любым
полем). Для $L \supseteq K$ множество $E(K)$ "--- подгруппа $E(L)$.
\end{remark}

\checkyourself{Чему равна сумма $P + Q + R$ трёх точек кривой, лежащих на одной невертикальной
прямой? А если прямая касается кривой в точке $P$ и проходит через $R$?}

\subsection{Точки порядка 2 и 3}

Обозначим $nP = P + \cdots + P$ ($n$ слагаемых), $0P = \ECpt$, $(-n)P = n(-P)$. \emph{Порядок}
точки $P$ "--- наименьшее $n \ge 1$ с $nP = \ECpt$.

\begin{proposition}\label{l8:order2}
Точка $P = (x_0, y_0) \ne \ECpt$ имеет порядок $2$ тогда и только тогда, когда $y_0 = 0$, т.\,е.
$x_0$ "--- корень $f$. Поэтому $E(K)$ содержит $0$, $1$ или $3$ точки порядка~$2$ (по числу корней
$f$ в $K$), а над $\overline K$ "--- ровно три, и
$E[2] = \{P : 2P = \ECpt\} \cong \ZZ/2 \times \ZZ/2$.
\end{proposition}

\begin{proof}
$2P = \ECpt \iff P = -P \iff y_0 = -y_0 \iff y_0 = 0$ (так как $2 \ne 0$ в $K$). Над $\overline K$
корни $f$ различны (теорема~\ref{l8:disc}), поэтому $E[2] = \{\ECpt, (e_1, 0), (e_2, 0), (e_3, 0)\}$
"--- группа из четырёх элементов, в которой каждый элемент сам себе обратен, т.\,е.
$\ZZ/2 \times \ZZ/2$.
\end{proof}

\begin{proposition}\label{l8:order3}
Точка $P = (x, y) \ne \ECpt$ имеет порядок $3$ тогда и только тогда, когда $x$ "--- корень
многочлена
\[
  \psi_3(x) = 3x^4 + 6ax^2 + 12bx - a^2.
\]
Геометрически это \emph{точки перегиба}: касательная в $P$ пересекает кривую в $P$
с кратностью $3$.
\end{proposition}

\begin{proof}
$3P = \ECpt \iff 2P = -P$. Точка порядка~$2$ не подходит ($2P = \ECpt \ne -P$), поэтому $y \ne 0$,
$\lambda = \frac{3x^2 + a}{2y}$, и $2P = -P$ равносильно $x(2P) = x$ (тогда $2P = \pm P$, а
$2P = P$ невозможно, так как $P \ne \ECpt$). Условие $\lambda^2 - 2x = x$, т.\,е.
$(3x^2 + a)^2 = 12x\,y^2 = 12x(x^3 + ax + b)$, после раскрытия скобок принимает вид
$9x^4 + 6ax^2 + a^2 = 12x^4 + 12ax^2 + 12bx$, т.\,е. $\psi_3(x) = 0$. В терминах леммы~\ref{l8:line}
$x(2P) = x$ означает, что третий корень многочлена $g$ для касательной тоже равен $x$.
\end{proof}

Над $\overline K$ многочлен $\psi_3$ имеет $4$ различных корня, каждый даёт две точки
$(x, \pm y)$, и вместе с $\ECpt$ получаем $9$ точек: $E[3] \cong \ZZ/3 \times \ZZ/3$ (см.
\cite[разд.~3.1]{washington}). Над $\RR$ у кривой ровно три точки перегиба, одна из которых "---
$\ECpt$.

\section{Эллиптические кривые над конечными полями}

\subsection{Группа $E(\GF{p})$}

Пусть $p > 3$ "--- простое, $a, b \in \GF{p}$, $4a^3 + 27b^2 \not\equiv 0 \pmod p$. Формулы
сложения те же, только все операции выполняются по модулю $p$, а деление "--- умножение на
обратный элемент (расширенный алгоритм Евклида, лекция~3 первого семестра). «Картинки» уже нет "---
точки кривой образуют конечное множество разбросанных точек плоскости $\GF{p}^2$, но группа
та же.

$E(\GF{p})$ "--- \emph{конечная} абелева группа: для каждого из $p$ значений $x$ имеется не
более двух $y$, поэтому $\#E(\GF{p}) \le 2p + 1$. По теореме Лагранжа порядок любой точки делит
$\#E(\GF{p})$.

\begin{example}\label{l8:ex-f13}
$E\colon y^2 = x^3 + 3x + 8$ над $\GF{13}$: $4 \cdot 27 + 27 \cdot 64 = 1836 \equiv 3 \not\equiv 0$.
Квадраты в $\GF{13}$: $\{1, 3, 4, 9, 10, 12\}$ ($1^2, 4^2, 2^2, 3^2, 6^2, 5^2$). Перебирая $x$:
\[
\begin{array}{c|ccccccccccccc}
x & 0 & 1 & 2 & 3 & 4 & 5 & 6 & 7 & 8 & 9 & 10 & 11 & 12\\ \hline
f(x) & 8 & 12 & 9 & 5 & 6 & 5 & 8 & 8 & 11 & 10 & 11 & 7 & 4\\
\legendre{f(x)}{13} & -1 & 1 & 1 & -1 & -1 & -1 & -1 & -1 & -1 & 1 & -1 & -1 & 1\\
y & & 5, 8 & 3, 10 & & & & & & & 6, 7 & & & 2, 11
\end{array}
\]
\[
  E(\GF{13}) = \{\ECpt, (1, 5), (1, 8), (2, 3), (2, 10), (9, 6), (9, 7), (12, 2), (12, 11)\},
  \qquad \#E(\GF{13}) = 9.
\]
Сложим $P = (9, 7)$ и $Q = (1, 8)$: $\lambda = \frac{8 - 7}{1 - 9} = \frac{1}{-8} \equiv 5^{-1}
\equiv 8$ (так как $5 \cdot 8 = 40 = 3 \cdot 13 + 1$), $x_3 = 64 - 9 - 1 = 54 \equiv 2$,
$y_3 = 8(9 - 2) - 7 = 49 \equiv 10$: $P + Q = (2, 10)$. Удвоим $P$:
$\lambda = \frac{3 \cdot 81 + 3}{14} \equiv \frac{12}{1} = 12$, $x_3 = 144 - 18 = 126 \equiv 9$,
$y_3 = 12 \cdot 0 - 7 \equiv 6$: $2P = (9, 6) = -P$, т.\,е. $P$ имеет порядок~$3$. Это
согласуется с предложением~\ref{l8:order3}: здесь $\psi_3(x) = 3x^4 + 18x^2 + 96x - 9$, и при
$x = 9$ имеем $9^2 \equiv 3$, $9^4 \equiv 9$,
$\psi_3(9) \equiv 3 \cdot 9 + 18 \cdot 3 + 96 \cdot 9 - 9 = 936 = 72 \cdot 13 \equiv 0$. Точка $Q = (1, 8)$ имеет
порядок $9$, так что группа циклическая: $E(\GF{13}) \cong \ZZ/9$.
\end{example}

\taskexample{ec-point-addition}

\checkyourself{Почему в примере~\ref{l8:ex-f13} порядок точки $(9, 7)$ мог быть только $1$, $3$
или $9$? Какие порядки возможны для точек кривой с $\#E = 10$?}

\subsection{Подсчёт точек и символ Лежандра}

Напомним (лекция~7), что \emph{символ Лежандра} $\legendre{c}{p}$ равен $0$ при
$p \mid c$, $1$, если $c$ "--- ненулевой квадрат по модулю $p$, и $-1$ иначе; по критерию Эйлера
$\legendre{c}{p} \equiv c^{(p-1)/2} \pmod p$.

\begin{lemma}\label{l8:sqcount}
Число решений сравнения $y^2 \equiv c \pmod p$ ($p > 2$) равно $1 + \legendre{c}{p}$.
\end{lemma}

\begin{proof}
При $c \equiv 0$ решение одно: $y = 0$ ($\GF{p}$ "--- поле, делителей нуля нет). Если $c$ "---
ненулевой квадрат, $c \equiv y_0^2$, то решения "--- корни многочлена
$y^2 - y_0^2 = (y - y_0)(y + y_0)$, т.\,е. ровно два различных числа $\pm y_0$ ($y_0 \ne -y_0$ при
$p > 2$). Если $c$ "--- не квадрат, решений нет.
\end{proof}

\begin{theorem}\label{l8:count}
Для эллиптической кривой $E\colon y^2 = f(x) = x^3 + ax + b$ над $\GF{p}$, $p > 3$,
\[
  \#E(\GF{p}) = p + 1 + \sum_{x \in \GF{p}} \legendre{f(x)}{p}.
\]
\end{theorem}

\begin{proof}
Аффинные точки с абсциссой $x$ "--- это решения $y^2 = f(x)$; по лемме~\ref{l8:sqcount} их
$1 + \legendre{f(x)}{p}$. Суммируя по всем $p$ значениям $x$ и добавляя $\ECpt$, получаем
$\#E(\GF{p}) = 1 + \sum_x \bigl(1 + \legendre{f(x)}{p}\bigr) = p + 1 + \sum_x \legendre{f(x)}{p}$.
\end{proof}

В примере~\ref{l8:ex-f13} сумма символов Лежандра равна $4 - 9 = -5$, и $\#E = 14 - 5 = 9$.

\begin{algorithm}[наивный подсчёт точек]\label{l8:countalg}
\begin{algorithmic}[1]
\Require простое $p > 3$, $a, b \in \GF{p}$ с $4a^3 + 27b^2 \ne 0$.
\Ensure $N = \#E(\GF{p})$.
\State $N \gets p + 1$
\For{$x = 0, 1, \dots, p - 1$}
  \State $c \gets x^3 + ax + b \bmod p$
  \State $N \gets N + \legendre{c}{p}$ \Comment{критерий Эйлера или закон взаимности}
\EndFor
\State \Return $N$
\end{algorithmic}
\end{algorithm}

Алгоритм выполняет $p$ вычислений символа Лежандра, т.\,е. $\bigO(p \log^2 p)$ битовых операций
(символ Якоби через закон взаимности считается за $\bigO(\log^2 p)$) "--- экспоненциально
по длине $\log p$ входа. Для криптографических $p \approx 2^{256}$ это немыслимо.
Р.~Схоф~\cite{schoof} в 1985 г. построил полиномиальный алгоритм подсчёта
($\bigO(\log^8 p)$ в исходном варианте); его улучшения (алгоритм SEA) считают точки на кривых
криптографических размеров за секунды. Для $p$ небольших учебных размеров удобно сначала выписать
все квадраты по модулю $p$, а затем для каждого $x$ проверять, есть ли $f(x)$ в этом списке.

\taskexample{ec-point-counting}

\subsection{Теорема Хассе}

Если бы значения $\legendre{f(x)}{p}$ были «случайными» $\pm 1$, сумма $p$ таких слагаемых
была бы порядка $\sqrt p$. Оказывается, это верно в точности, и с явной константой.

\begin{theorem}[Хассе]\label{l8:hasse}
Для эллиптической кривой $E$ над $\GF{p}$
\[
  \abs{\#E(\GF{p}) - (p + 1)} \le 2\sqrt p,
\]
т.\,е. $p + 1 - 2\sqrt p \le \#E(\GF{p}) \le p + 1 + 2\sqrt p$. (Без доказательства; см.
\cite[разд.~4.2]{washington}, \cite[гл.~V, \S\,1]{silverman-aec}.)
\end{theorem}

Число $t = p + 1 - \#E(\GF{p})$ называется \emph{следом Фробениуса}; теорема Хассе утверждает,
что $t^2 \le 4p$, а по теореме~\ref{l8:count} $t = -\sum_x \legendre{f(x)}{p}$. Интервал
$[p + 1 - 2\sqrt p,\ p + 1 + 2\sqrt p]$ называют \emph{интервалом Хассе}. В примере~\ref{l8:ex-f13}
$t = 5$, $2\sqrt{13} \approx 7{,}21$. Теорема М.~Дойринга (Deuring) дополняет оценку Хассе:
для каждого целого $t$ с $\abs{t} \le 2\sqrt p$ существует кривая над $\GF{p}$ с
$\#E = p + 1 - t$~\cite[разд.~7.3]{crandall-pomerance}.

Следствия для криптографии и алгоритмов:
\begin{itemize}
  \item $\#E(\GF{p}) \approx p$, т.\,е. группа точек имеет «тот же размер», что и $\Znx{p}$:
    кривая над $256$-битным полем даёт группу порядка около $2^{256}$. Для протоколов выбирают
    кривые, у которых $\#E = h \cdot q$ с большим простым $q$ и малым \emph{кофактором} $h$
    (часто $h = 1$), иначе DLP сводится к малым подгруппам (метод Полига"--~Хеллмана, лекция~2).
  \item Порядки групп кривых над $\GF{p}$ «пробегают» весь интервал Хассе, когда кривая меняется;
    это используется в методе Ленстры (лекция~10): среди случайных кривых найдутся такие, у
    которых $\#E(\GF{p})$ "--- гладкое число.
  \item Зная $\#E(\GF{p})$ с точностью до интервала длины $4\sqrt p$, порядок точки можно найти
    алгоритмом «шаг младенца "--- шаг великана» за $\bigO(p^{1/4})$ операций.
\end{itemize}

\begin{remark}[кручение]
Пусть $d \in \GF{p}^*$ "--- не квадрат. Кривая $E'\colon y^2 = \tilde f(x) = x^3 + d^2 a x + d^3 b$
называется \emph{квадратичным кручением} $E$. Так как $\tilde f(dx) = d^3 f(x)$, а $dx$ пробегает
$\GF{p}$ вместе с $x$,
\[
  \sum_x \legendre{\tilde f(x)}{p} = \sum_x \legendre{\tilde f(dx)}{p}
  = \legendre{d}{p}^3 \sum_x \legendre{f(x)}{p} = -\sum_x \legendre{f(x)}{p},
\]
и по теореме~\ref{l8:count} $\#E(\GF{p}) + \#E'(\GF{p}) = 2p + 2$: у кручения кривой из примера~\ref{l8:ex-f13} ровно
$28 - 9 = 19$ точек.
\end{remark}

\checkyourself{Какие значения может принимать $\#E(\GF{11})$? А $\#E(\GF{101})$?}

\subsection{Строение группы $E(\GF{p})$}

Группа $E(\GF{p})$ не обязана быть циклической.

\begin{theorem}\label{l8:structure}
Для эллиптической кривой $E$ над $\GF{p}$ существуют натуральные $n_1, n_2$, такие что
\[
  E(\GF{p}) \cong \ZZ/n_1 \times \ZZ/n_2, \qquad n_1 \mid n_2, \qquad n_1 \mid p - 1.
\]
(Без доказательства; см.~\cite[разд.~4.1]{washington}.)
\end{theorem}

Первая часть следует из теоремы о строении конечных абелевых групп и того, что подгруппа
$n$-кручения $E[n]$ над $\overline{\GF{p}}$ при $p \nmid n$ изоморфна $\ZZ/n \times \ZZ/n$
(так что в разложении не больше двух множителей); условие $n_1 \mid p - 1$ "--- следствие
свойств спаривания Вейля.

\begin{example}
$E\colon y^2 = x^3 - x$ над $\GF{7}$. Все три корня $0, 1, -1$ многочлена $x^3 - x$ лежат в
$\GF{7}$, поэтому по предложению~\ref{l8:order2} у $E(\GF{7})$ три точки порядка~$2$: $(0, 0)$,
$(1, 0)$, $(6, 0)$, и группа не циклическая (в циклической группе элемент порядка~$2$ один).
Подсчёт даёт ещё $(4, 2)$, $(4, 5)$, $(5, 1)$, $(5, 6)$; $\#E(\GF{7}) = 8$, и
$E(\GF{7}) \cong \ZZ/2 \times \ZZ/4$ (здесь $n_1 = 2 \mid 6$).
\end{example}

\taskexample{ec-point-order}

\checkyourself{Может ли быть $E(\GF{5}) \cong \ZZ/3 \times \ZZ/3$? Почему группа $E(\GF{p})$,
порядок которой свободен от квадратов, циклическая?}

\section{Изоморфизмы кривых и $j$-инвариант}

Замена $x = u^2 x'$, $y = u^3 y'$ ($u \in K^*$) переводит $y^2 = x^3 + ax + b$ в
$y'^2 = x'^3 + u^{-4} a x' + u^{-6} b$ (поделите уравнение на $u^6$). Она сохраняет прямые
вида $y = \lambda x + \nu$ и вертикальные прямые, а потому и закон сложения: отображение
$(x, y) \mapsto (u^{-2}x, u^{-3}y)$, $\ECpt \mapsto \ECpt$, "--- изоморфизм групп. Можно
показать, что других изоморфизмов между кривыми в короткой форме (переводящих $\ECpt$ в $\ECpt$)
нет~\cite[разд.~2.7]{washington}.

\begin{definition}
\emph{$j$-инвариантом} кривой $y^2 = x^3 + ax + b$ называется
\[
  j(E) = 1728 \cdot \frac{4a^3}{4a^3 + 27b^2}.
\]
\end{definition}

При замене $(a, b) \mapsto (u^{-4} a, u^{-6} b)$ числитель и знаменатель умножаются на
$u^{-12}$, поэтому \emph{изоморфные кривые имеют одинаковый $j$-инвариант}. Верно и обратное
над алгебраически замкнутым полем: кривые с одинаковым $j$ изоморфны над $\overline K$ (но,
вообще говоря, не над $K$: так, кривая и её квадратичное кручение имеют одинаковый $j$, но
различное число точек). Например, $j = 0$ у всех кривых $y^2 = x^3 + b$ (к ним относится
secp256k1: $y^2 = x^3 + 7$ над $\GF{p}$, $p = 2^{256} - 2^{32} - 977$), $j = 1728$ у всех кривых
$y^2 = x^3 + ax$; для кривой примера~\ref{l8:ex-q} $j = \frac{54000}{11}$.

\checkyourself{Докажите, что для любого $j_0 \in \GF{p}$ существует эллиптическая кривая над
$\GF{p}$ с $j(E) = j_0$. (Указание: при $j_0 \ne 0, 1728$ возьмите $a = b$.)}

\section{Вычисления без обращений}

Каждое сложение по алгоритму~\ref{l8:addalg} требует одного обращения в $\GF{p}$, а оно
в десятки раз дороже умножения. В \emph{проективных координатах} точку представляют тройкой
$(X : Y : Z)$ и выполняют сложение только умножениями, откладывая единственное обращение до
конца вычислений. Особенно удобны \emph{координаты Якоби}: $(X : Y : Z)$ соответствует аффинной
точке $(X/Z^2, Y/Z^3)$, уравнение кривой принимает вид $Y^2 = X^3 + aXZ^4 + bZ^6$, а
$\ECpt = (1 : 1 : 0)$. Формулы удвоения:
\[
  S = 4XY^2, \quad M = 3X^2 + aZ^4, \qquad
  X' = M^2 - 2S, \quad Y' = M(S - X') - 8Y^4, \quad Z' = 2YZ.
\]
Это около десяти умножений и ни одного обращения (проверьте подстановкой, что
$(X'/Z'^2, Y'/Z'^3)$ совпадает с аффинным удвоением точки $(X/Z^2, Y/Z^3)$; мы проверили это
в Sage). Аналогичные формулы есть и для сложения.

Быстрое вычисление кратных $nP$ (лекция~9) сводится к $\bigO(\log n)$ удвоениям и сложениям,
и там выигрыш от проективных координат особенно заметен. В методе Ленстры (лекция~10), наоборот,
\emph{невозможность} обратить знаменатель по составному модулю и есть цель вычислений.

\section{Кривые над полями характеристики 2 и 3}\label{l8:char23}

При $\operatorname{char} K = 2$ или $3$ привести общее уравнение~\eqref{l8:general} к короткой
форме нельзя (требуются деления на $2$ и $3$). При $\operatorname{char} K = 2$ кривая
$y^2 = x^3 + ax + b$ всегда особая: $F_y = 2y \equiv 0$, а $F_x = x^2 + a$ обращается в нуль при
$x_0 = \sqrt a$, так что $(x_0, \sqrt{f(x_0)})$ "--- особая точка над $\overline K$. При
$\operatorname{char} K = 3$ от слагаемого $x^2$ нельзя избавиться, и короткая форма описывает
лишь часть кривых. Используют другие нормальные формы:
\begin{itemize}
  \item $\operatorname{char} K = 2$: $y^2 + xy = x^3 + a_2 x^2 + a_6$, $a_6 \ne 0$ (обыкновенные
    кривые) или $y^2 + a_3 y = x^3 + a_4 x + a_6$, $a_3 \ne 0$ (суперсингулярные);
  \item $\operatorname{char} K = 3$: $y^2 = x^3 + a_2 x^2 + a_4 x + a_6$.
\end{itemize}
Правило хорд и касательных и все свойства группы сохраняются; меняются лишь формулы. Например,
в общей форме~\eqref{l8:general} $-(x, y) = (x, -y - a_1 x - a_3)$, т.\,е. на кривой
$y^2 + xy = x^3 + a_2 x^2 + a_6$ над $\GF{2^m}$ имеем $-(x, y) = (x, y + x)$. Кривые над
$\GF{2^m}$ привлекательны для аппаратной реализации и встречались в стандартах, но в
современной практике преобладают кривые над $\GF{p}$.

\checkyourself{Выведите формулу $-(x, y) = (x, -y - a_1 x - a_3)$ для кривой в общей
форме~\eqref{l8:general}: какие две точки кривой лежат на вертикальной прямой $x = x_0$?}

\checkyourselfsection

\begin{furtherreading}
  \readitem[разд.~6.1--6.2]{hps}{эллиптические кривые над $\RR$ и над $\GF{p}$, закон сложения,
    теорема Хассе "--- основной текст лекции}
  \readitem[гл.~7, разд.~7.1--7.3, 7.5]{crandall-pomerance}{арифметика кривых, проективные
    координаты, теоремы Хассе и Дойринга, подсчёт точек (алгоритм Схофа)}
  \readitem[гл.~VI]{koblitz}{эллиптические кривые и криптосистемы на них}
  \readitem[гл.~2, 4]{washington}{уравнения Вейерштрасса, доказательство ассоциативности
    (разд.~2.4), строение $E(\GF{q})$ и теорема Хассе (гл.~4)}
  \readitem[гл.~III, V]{silverman-aec}{строгое изложение: закон сложения, изоморфизмы,
    $j$-инвариант, кривые над конечными полями}
  \readitem{schoof}{полиномиальный алгоритм подсчёта точек}
  \readitem{koblitz-ecc}{одна из двух статей, с которых началась эллиптическая криптография}
  \readitem[лекции~1--2]{mit-18783}{конспекты курса Э.~Сазерленда: закон сложения,
    проективные координаты}
\end{furtherreading}
\end{lecture}
